\documentclass[11pt]{article}
\usepackage[margin=1in]{geometry}
\usepackage{amsmath,amssymb,amsthm}
\usepackage{graphicx}
\usepackage{booktabs}
\usepackage{enumitem}
\usepackage[hidelinks]{hyperref}

\newtheorem{theorem}{Theorem}
\newtheorem{corollary}[theorem]{Corollary}
\newtheorem{definition}[theorem]{Definition}
\newcommand{\E}{E^1}
\newcommand{\D}{\mathbb D}
\newcommand{\T}{\mathbb T}

\title{A Three-Sheet Inverse-Monodromy Criterion for the\\
Quadratic Smirnov Uniqueness Condition}
\author{Candidate partial-result packet for arXiv:2601.04146}
\date{9 August 2026}

\begin{document}
\maketitle

\begin{center}
\fbox{\textbf{Status: substantial partial result, likely valid; pending specialist review}}
\end{center}

\section{Source problem}

Fricain, Grivaux, Ostermann, and Yakubovich study commutants and
embeddability of Toeplitz operators with smooth symbols in
\cite{FGOYembedding}.  If $\Omega$ is a component of
$\sigma(T_F)\setminus F(\T)$ with $|\operatorname{wind}_F(\Omega)|=2$, put
\[
 \Gamma=\partial\Omega\cap\partial\Omega_1,
 \qquad \zeta=1/F^{-1}
\]
on the regular arcs of the parametrized curve.  Condition (1) of their
Theorem~6.7 is
\begin{equation}\label{eq:source-condition}
 u+\zeta v+\zeta^2w=0\quad\text{a.e. on }\Gamma,
 \qquad u,v,w\in E^1(\Omega)
 \quad\Longrightarrow\quad u=v=w=0.
\end{equation}
Immediately after the theorem, the authors state that they do not know a
geometric condition implying \eqref{eq:source-condition}.  Their earlier
Property (P) treats the corresponding two-term condition
\cite{FGOhypercyclicity}, but does not supply the three-term criterion.

\begin{figure}[ht]
  \centering
  \includegraphics[width=\textwidth]{figures/open_problem_crop.png}
  \caption{Theorem 6.7 and the stated missing geometric condition in
  \cite[p.~30]{FGOYembedding}.}
\end{figure}

\section{The inverse-cover condition}

The relevant geometric datum is not merely the planar set $F(\T)$ but the
sheet structure of the parametrized reciprocal inverse.  This distinction is
unavoidable: Example~6.6 of the source gives the same planar curve square-root
and cube-root parametrizations with opposite commutant behavior.

\begin{definition}\label{def:cover}
Let $\Omega$ be a bounded Smirnov domain with piecewise analytic boundary and
let $\Gamma\subset\partial\Omega$ contain an analytic arc of positive length.
We say that $(\Omega,\zeta)$ has a \emph{bounded finite inverse continuation
of sheet degree $d$} if there are
\begin{enumerate}[label=(\roman*)]
\item a connected Riemann surface $X$ and a finite proper holomorphic map
      $\pi:X\to\Omega$ of degree $d$;
\item a function $\Psi\in H^\infty(X)$ which takes $d$ distinct values on a
      generic fiber of $\pi$;
\item a positive-length analytic subarc $\gamma\subset\Gamma$ and a
      one-sided local section $s$ of $\pi$ over an interior collar of
      $\gamma$ such that the non-tangential boundary values of
      $\Psi\circ s$ equal $\zeta$ a.e. on $\gamma$.
\end{enumerate}
The continuation has \emph{full boundary compatibility} if, on almost every
regular subarc of $\Gamma$, $\zeta$ is the boundary value of some local sheet
of $\Psi$.
\end{definition}

Thus $d$ counts distinct analytic values, not merely points in an artificially
enlarged covering.  For $\zeta(\lambda)=\lambda^{1/d}$, take
$X=\D$, $\pi(x)=x^d$, and $\Psi(x)=x$.

\section{Proof intuition}

The boundary relation is first converted, by Privalov uniqueness on a small
one-sided collar, into a holomorphic identity on one inverse sheet.  Lifting
to $X$ and using connectedness propagates that identity to every sheet.  Over
a regular value $\lambda$, the same polynomial in one variable must therefore
vanish at all $d$ distinct values of $\Psi$ in the fiber.  A polynomial of
degree at most $r<d$ cannot do so unless all its coefficients vanish.

The converse is just as rigid.  If $d\le r$, multiply the linear factors
$t-\Psi(x)$ over a fiber.  The elementary symmetric coefficients descend to
bounded holomorphic functions on $\Omega$, including across branch values,
and the resulting nonzero degree-$d$ polynomial annihilates every boundary
sheet.  This makes the sheet threshold sharp.

\section{Finite-sheet classification}

\begin{theorem}\label{thm:classification}
Let $r\ge0$, and suppose $(\Omega,\zeta)$ has a bounded finite inverse
continuation of sheet degree $d$ in the sense of Definition~\ref{def:cover}.
If $d>r$, then
\begin{equation}\label{eq:relation-r}
 \sum_{k=0}^{r}a_k\zeta^k=0\quad\text{a.e. on }\Gamma,
 \qquad a_0,\ldots,a_r\in E^1(\Omega),
\end{equation}
implies $a_0=\cdots=a_r=0$ on $\Omega$.

If full boundary compatibility holds and $d\le r$, the uniqueness conclusion
fails.  Consequently, within the fully compatible bounded finite-continuation
class, degree-$r$ boundary uniqueness holds if and only if $d>r$.
\end{theorem}

\begin{proof}
Put
\[
 P(\lambda,t)=\sum_{k=0}^{r}a_k(\lambda)t^k.
\]
Shorten the seed arc $\gamma$ from Definition~\ref{def:cover} if necessary,
and choose a piecewise analytic one-sided collar $W\subset\Omega$ whose
artificial boundary is compactly contained in $\Omega$.  Let
$\psi=\Psi\circ s$ on $W$.  It is bounded.

Every $a_k\in E^1(\Omega)$ belongs to the Nevanlinna class
$\mathcal N(\Omega)$.  Restricting a quotient of bounded holomorphic
functions shows that $a_k|_W\in\mathcal N(W)$.  Since $\psi$ is bounded and
the Nevanlinna class is a field,
$P(\lambda,\psi(\lambda))\in\mathcal N(W)$.  Relation
\eqref{eq:relation-r} makes its non-tangential boundary values zero on a
positive-length arc, so the standard Privalov boundary uniqueness theorem
for Nevanlinna-class functions gives
\begin{equation}\label{eq:seed-identity}
 P(\lambda,\psi(\lambda))=0\qquad(\lambda\in W).
\end{equation}

The function
\[
 x\longmapsto P(\pi(x),\Psi(x))
\]
is holomorphic on $X$.  Equation \eqref{eq:seed-identity} says that it
vanishes on an open subset.  Since $X$ is connected, it vanishes identically.
For every regular value $\lambda$ at which $\Psi$ separates the fiber, the
polynomial $P(\lambda,t)$ therefore has the $d$ distinct roots
\[
 \{\Psi(x):x\in\pi^{-1}(\lambda)\}.
\]
If $d>r$, this polynomial is identically zero.  Thus every $a_k$ vanishes on
the nonempty open set of such regular values, and the identity theorem gives
$a_k=0$ throughout $\Omega$.

For the converse, initially at regular values define
\[
 Q(\lambda,t)=\prod_{x\in\pi^{-1}(\lambda)}(t-\Psi(x))
 =t^d+c_{d-1}(\lambda)t^{d-1}+\cdots+c_0(\lambda).
\]
The coefficients are symmetric in the fiber, hence are single-valued and
holomorphic away from the branch values.  If $\|\Psi\|_\infty\le M$, then
$|c_{d-j}|\le \binom{d}{j}M^j$.  The removable-singularity theorem extends
each coefficient through the discrete branch-value set, so
$c_j\in H^\infty(\Omega)\subset E^1(\Omega)$.  Full boundary compatibility
implies
\[
 Q(\lambda,\zeta(\lambda))=0\quad\text{a.e. on }\Gamma.
\]
This is a nonzero relation of degree $d\le r$; append zero higher
coefficients when $d<r$.  Uniqueness therefore fails.
\end{proof}

\begin{corollary}[Sharp three-sheet criterion]\label{cor:quadratic}
If $(\Omega,\zeta)$ has a bounded inverse continuation with at least three
sheets, then the source condition \eqref{eq:source-condition} holds.  Under
full boundary compatibility in the finite-sheet class, it holds exactly when
$d\ge3$.
\end{corollary}

\begin{proof}
Apply Theorem~\ref{thm:classification} with $r=2$.  For $d=2$, the explicit
obstruction is
\[
 \zeta^2-(\psi_1+\psi_2)\zeta+\psi_1\psi_2=0,
\]
where the symmetric sum and product descend to $H^\infty(\Omega)$.  The
one-sheet case already has a nonzero linear relation.
\end{proof}

Combining Corollary~\ref{cor:quadratic} with Theorem~6.7 of
\cite{FGOYembedding} shows that, when no component of winding modulus greater
than two occurs, the three-sheet condition forces the commutant of the model
operator to consist only of multipliers and yields the paper's embeddability
criterion.

\section{An explicit infinite family}

For each integer $m\ge3$, let $a_m=2\pi(m-1)/m$ and define
\begin{equation}\label{eq:Fm}
 F_m(e^{i\theta})=
 \begin{cases}
 2-m+(m-1)e^{-im\theta/(m-1)},&0\le\theta\le a_m,\\
 e^{-im\theta},&a_m\le\theta\le2\pi.
 \end{cases}
\end{equation}
At $\theta=0,a_m,2\pi$, the pieces have the same value $1$ and matching
first derivative $-im$.  Both pieces have constant nonzero speed $m$; the
derivative is globally Lipschitz.  Thus $F_m$ is $C^{1,1}$ as a function on
$\T$ and satisfies the source's smoothness hypothesis (H1) for every required
H\"older exponent below one.

The image consists of two clockwise circles: the circle centered at $2-m$ of
radius $m-1$, and the unit circle.  They are internally tangent only at $1$.
Consequently (H2) holds, all winding numbers are nonpositive, and
\[
 \operatorname{wind}_{F_m}(\lambda)=
 \begin{cases}
 -2,&|\lambda|<1,\\
 -1,&|\lambda-(2-m)|<m-1\text{ and }|\lambda|>1,\\
 0,&\text{outside the outer circle}.
 \end{cases}
\]
On the boundary of the doubly wound component $\Omega=\D$,
\[
 \zeta(\lambda)=
 \exp\!\left(\frac1m\operatorname{Log}_{(0,2\pi)}\lambda\right).
\]
It has the continuation model
\[
 X=\D,\qquad \pi(x)=x^m,\qquad\Psi(x)=x,
\]
of sheet degree $m$.  Corollary~\ref{cor:quadratic} proves
\eqref{eq:source-condition} for every $m\ge3$.  The member $m=3$ is the
cube-root example already proved directly in Example~6.6 of the source; the
classification supplies the structural explanation and all $m\ge4$ at once.

\section{Verification and limitations}

The current verdict is \textbf{likely valid} (confidence $91/100$).

\begin{center}
\begin{tabular}{p{0.27\textwidth}p{0.64\textwidth}}
\toprule
Check & Finding\\
\midrule
Boundary-to-interior step & Since $E^1(\Omega)\subset\mathcal N(\Omega)$,
restriction to the collar and multiplication by a bounded inverse sheet stay
in $\mathcal N(W)$; Privalov boundary uniqueness applies.\\
Continuation & The relation is lifted to one connected surface; no comparison
between unrelated local sectors is made.\\
Branch values & Symmetric coefficients are bounded near the discrete branch
set and extend removably.\\
Sharpness & The fiber polynomial gives an explicit nonzero relation whenever
$d\le r$.\\
Family sanity check & \texttt{code/verify\_family.py} checks junction values,
first derivatives, representative winding numbers, and the three-root
Vandermonde determinant for $m=3,\ldots,20$.  This is not a proof.\\
\bottomrule
\end{tabular}
\end{center}

The theorem requires bounded finite analytic continuation of the reciprocal
inverse.  It does not classify arbitrary symbols whose inverse continuation
has infinitely many sheets, natural boundaries, poles, or no continuation
beyond a collar.  It is also a condition on the sheet geometry of the
\emph{parametrized} curve, not on the unparametrized planar image alone.  The
source's square-root/cube-root pair shows why the latter distinction is
essential.

For these reasons the packet is conservatively classified as a substantial
partial result.  A specialist should decide whether a finite inverse-cover
condition counts as the geometric condition requested by the source.  Within
the bounded finite-continuation class, Theorem~\ref{thm:classification} is a
full sharp classification.

\section{Bounded novelty search}

On 9 August 2026, the run's lightweight registry, solution, attempt, and
proof-gap indexes were searched for arXiv:2601.04146, the displayed three-term
relation, inverse monodromy, and finite-sheet variants.  The only relevant
prior attempt, on arXiv:2502.03303, concerns different hypercyclicity
questions.  Official-arXiv web searches used the source title together with
``monodromy,'' ``reciprocal inverse monodromy commutant,'' ``Smirnov quadratic
relation,'' and the displayed relation.  They found the source paper but no
paper stating this criterion or classification.  This bounded search supports
but does not guarantee novelty.

\small
\begin{thebibliography}{9}

\bibitem{FGOYembedding}
E. Fricain, S. Grivaux, M. Ostermann, and D. Yakubovich,
\emph{Embedding of Toeplitz operators with smooth symbols into strongly
continuous semigroups}, arXiv:2601.04146 (2026), especially Example 6.6 and
Theorem 6.7.

\bibitem{FGOhypercyclicity}
E. Fricain, S. Grivaux, and M. Ostermann,
\emph{Hypercyclicity of Toeplitz operators with smooth symbols},
arXiv:2502.03303 (2025), especially Property (P) and Proposition 4.9.

\end{thebibliography}
\normalsize

\end{document}
